\section{Note on Sources and Delivery}

For an adult parish assembly at the universal 1962 Twenty-fourth and Last
Sunday after Pentecost, 22 November 2026. The speech develops the reviewed
interpretation ``The confident petitioner asks to be changed,'' with the
compatible distinction between present rescue and continuing danger from
``Deliverance begun, and the people who still wait.'' Its connections
between Matthew's warning, Paul's petition and the three orations are the
homily's own synthesis, not a whole-Mass argument attributed to a Father.
Bede expounds the immediate context of Mk 11:24, particularly verse 25,
rather than a separately quoted verse-24 lemma. Biblical exposition concerns
Mt 24:1--3, 15--35; Col 1:9--14; and Mk 11:22--26. The English biblical
context and short Epistle quotation follow the Douay--Rheims--Challoner.

The speech contains 1,207 words (TeXcount, excluding title and apparatus).
At an unhurried 110--120 words per minute, the estimated duration is
10.1--11.0 minutes. The prose was reread silently for sense, sentence length
and oral transitions; no audible rehearsal or timed human delivery is
claimed. The References and this note are outside the speech.

\section{References}

\begin{itemize}
\item \textbf{Appointed texts.}
\work{Missale Romanum}, Vatican typical edition, 1962,
pp. 422--424, nos. 1769--1778, especially Collect, Secret,
Communion and Postcommunion.
\url{https://media.churchmusicassociation.org/pdf/missale62.pdf}.
Francis Xavier Lasance and Francis Augustine Walsh,
\work{The New Roman Missal}, revised-1945 expression in its later
printing, Last Sunday after Pentecost, pp. 736--740.
The short Communion quotation follows this historical English,
not an approved English edition of the 1962 Missal.

\item \textbf{Scripture.} Douay--Rheims--Challoner,
registered Project Gutenberg no. 1581 per-book texts:
Mt 24:1--3, 15--35; Col 1:9--14; Mk 11:22--26.
The short Epistle quotation follows this English.
\url{https://www.gutenberg.org/ebooks/1581}.
Inherited biblical and liturgical wording retains its separate
rights status outside the project's licence.

\item \textbf{John Chrysostom.}
\work{Homilies on Matthew} 76, especially §2,
historical NPNF first-series vol. X English, registered CCEL text,
on Judean flight and the distinct final coming.
\work{Homilies on Colossians} 2, on Col 1:9--14,
historical English attributed to J. Ashworth, NPNF XIII,
registered CCEL text, on continued prayer and fitness for the gift.

\item \textbf{Thomas Aquinas.}
\work{Super Colossenses} 1, lectio 3, native anchor 87848:
prayer's ascent to God, fruitful knowledge, and being made worthy.
\url{https://www.corpusthomisticum.org/ccl.html}.
Original exposition of the checked Latin; no restricted modern
translation or delivery text is reproduced.

\item \textbf{Bede.}
\work{In Marci evangelium expositio} III, on Mk 11:25,
paragraph beginning \textit{Et cum stabitis ad orandum},
registered Migne-derived Latin transcription:
faith working by love, petition for pardon and forgiveness of others.

\item \textbf{Augustine.}
\work{De Trinitate} III.4.10, historical English
trans. Arthur West Haddan (1887), registered Book III text:
sacramental food and drink for spiritual health.
\end{itemize}
